The strategy executes in four stages --- threshold-wise tree training,
rule extraction, membership-function construction, and consequent
tuning --- wrapped by the nested protocol of
Section~\ref{sec:nested}. Algorithm~\ref{alg:method} summarizes the
complete procedure.

\begin{algorithm}[htbp]
  \caption{Tree-rule fuzzy estimation of the corrosion rate}
  \label{alg:method}
  \KwIn{records $(\mathbf{x}_i, v_i)$, normative limits $L = \{2, 8\}$
    [mpy], saturation level $v_{\max} = 20$ [mpy]}
  \KwOut{fuzzy estimator $\hat{v}(\mathbf{x})$; out-of-fold metrics}
  saturate the target: $v_i \leftarrow \min(v_i, v_{\max})$\;
  \ForEach{outer fold $(\mathcal{T}, \mathcal{E})$ of the stratified
    5-fold split}{
    \tcp{all selection inside the training fold $\mathcal{T}$}
    select depth, leaf size, pruning and leaf threshold of each tree
    by inner CV maximizing mean recall\;
    fit binomial trees: $T_2$ for $v > 2$, $T_8$ for $v > 8$ on
    $\mathcal{T}$\;
    collapse sibling leaves of equal class; extract root-to-leaf
    paths as rules: $T_2^{-} \to$ low, $T_2^{+} \to$ mid,
    $T_8^{+} \to$ high\;
    merge near-duplicate thresholds; simplify each rule to one
    interval per variable\;
    select by inner CV (MAE) the numeric-layer configuration:
    rule-source trees (recall-selected or deeper variants),
    transition-width factor, and consequent regularization\;
    \ForEach{variable used by the rules}{
      partition its range at the tree thresholds\;
      build one interval membership per region: sigmoid shoulders
      crossing $0.5$ at the thresholds, width from the distance
      threshold--median of the training data, scaled by the selected
      factor\;
    }
    initialize each consequent $c_r$ at the median of its class;
    adjust all $c_r$ by regularized least squares on the normalized
    activations, bounded to the normative range of each rule\;
    predict $\mathcal{E}$ with Eq.~(\ref{eq:fis}), applying to the
    mid rules the ordinal exclusion
    $w \leftarrow \min(w,\, 1 - \max w_{\mathrm{high}})$\;
  }
  report metrics over the union of the outer test folds\;
  refit the complete chain on all records for deployment (artifacts
  only)\;
\end{algorithm}


\subsection{Binomial trees at the normative limits}

One binomial decision tree per normative limit provides the
classification backbone: $T_2$ answers whether the rate exceeds
2 [mpy] and $T_8$ whether it exceeds 8 [mpy]. Both use balanced class
weights. The hyperparameters --- depth (3--5), minimum leaf size
(3--10), cost-complexity pruning, and the leaf probability threshold
$\tau \in \{0.4, 0.5, 0.6\}$ --- compete in the inner folds under the
criterion mean recall $+\ 0.5 \times$ recall ratio, which targets the
acceptance condition of the application instead of accuracy.

\subsection{Rule extraction}

Each root-to-leaf path is a conjunction of threshold conditions.
Sibling leaves of the same class merge recursively, because their
split does not change the decision; the diagrams of the deployed trees
therefore match the rulebase one-to-one. The extraction assigns the
negative paths of $T_2$ to the consequent \emph{low}, the positive
paths of $T_2$ to \emph{mid}, and the positive paths of $T_8$ to
\emph{high}. Thresholds of one variable that differ by less than 1\%
of its robust range (1st--99th percentile) merge into one value, and
the conditions of each rule simplify to a single interval per
variable. Because the positive region of $T_2$ contains the severe
region, every \emph{mid} rule carries the ordinal exclusion \emph{AND
NOT (severe conditions)}, implemented as
$w \leftarrow \min(w, 1 - \max w_{\mathrm{high}})$.

\subsection{Interval membership functions}

For each variable used by the rulebase, its sorted tree thresholds
$t_1 < \dots < t_k$ partition the observed range into $k+1$ intervals,
and each interval defines one linguistic term. The term membership is
the minimum of two logistic shoulders: a rising shoulder anchored at
the lower threshold and a falling shoulder anchored at the upper
threshold. Each shoulder crosses 0.5 exactly at its threshold, which
preserves the correspondence between the tree partition and the fuzzy
partition. The transition width equals the distance from the threshold
to the median of the training data inside the interval, multiplied by
a scale factor selected in the inner folds
($s \in \{0.03, 0.1\}$); widths remain bounded by the robust span
of the variable so that outliers cannot distort the shape. A tree
condition $x \leq t_j$ maps to the maximum membership of the terms
left of $t_j$, and $x > t_j$ to the maximum of the terms right of it.

\subsection{Consequent tuning and inference}

Each rule receives a singleton consequent $c_r$ [mpy], initialized at
the median of the training records of its class and then adjusted by
box-constrained least squares on the normalized activation matrix:
$c_r$ must remain inside the normative range of its class (low:
$[0,2]$; mid: $[2,8]$; high: $[8, v_{\max}]$), with a regularization
$\lambda \in \{0.3, 1.0\}$, selected in the inner folds, that pulls
each consequent toward its class median. The constraint replaces the
manual support adjustment of classical fuzzy design with an auditable,
norm-consistent optimization: the expert reviews the tuned supports
knowing that no rule can speak outside its severity range. For the
numeric layer, the inner loop also chooses the rule-source trees among
the recall-selected ones and two deeper variants (depth 6--7): finer
partitions can carry additional regression signal, while the reported
classification always comes from the recall-selected trees. Inference
follows Eq.~(\ref{eq:fis}); when no rule fires, the system returns
the training median.

\subsection{Benchmarks}

Three regression methods suited to reduced databases compete on the
identical outer folds: linear regression, a decision-tree regressor
(depth 3--6, leaf size 3--10, selected in inner folds), and support
vector regression ($C \in \{1, 10, 100\}$,
$\epsilon \in \{0.1, 1.0\}$). Ensemble regressors with large parameter
counts remain outside the comparison because the reduced size of the
database raises their overfitting risk and their predictions do not
admit rule-level audit. Every benchmark uses median imputation and
standardization fitted inside the training folds.
